%%%%%%%%%%%%%%%%%%%%%%%%%%%%%%%%%%%%%%%%%%%%%%%%%%%%%%%%%%%%%%%%%%%%%%
% Lớp 11 | Hàm số lượng giác và phương trình lượng giác | [Giảm] Phương trình lượng giác thường gặp
% Dạng 1D1?6-1: Phương trình bậc n theo một hàm số lượng giác
% Số câu: 3  (sinh tự động bởi tach_cau.py; từ nay thầy sửa trực tiếp file này)
%%%%%%%%%%%%%%%%%%%%%%%%%%%%%%%%%%%%%%%%%%%%%%%%%%%%%%%%%%%%%%%%%%%%%%

%%%%%%%%%%%%%%%% TRẮC NGHIỆM 4 PHƯƠNG ÁN %%%%%%%%%%%%%%%%

% nguồn: Ok_TracNghiemLop11_Lan1
\begin{ex}%[1D1H6-1]%[ma:b68762ed]
		 Giải phương trình $\cos 2x+2\cos x-3=0$.
	 \choice
			{$x=\pi+k2\pi,k\in\mathbb{Z}$} 
			 {\True $x=k2\pi,k\in\mathbb{Z}$}
			{$x=-\dfrac{\pi}{2}+k2\pi,k\in\mathbb{Z}$} 
			{$x=\dfrac{\pi}{2}+k2\pi,k\in\mathbb{Z}$} 
	\loigiai{ 
			$\cos 2x+2\cos x-3=0\Leftrightarrow 2\cos^2x+2\cos x-4=0\Leftrightarrow \hoac{&\cos x=1\\&  \cos x=-2\;\text{(loại)}}\Leftrightarrow x=k2\pi,k\in\mathbb{Z}$.
		 }
\end{ex}


%%%%%%%%%%%%%%%% ĐÚNG/SAI %%%%%%%%%%%%%%%%

% nguồn: Ok_dungsai_lop11_Lan1
\begin{ex}%[1D1H6-1]%[ma:66dbdb13]
 Phương trình $4\sin^4x+12\cos^2x-7=0$ tương đương với phương trình nào dưới  đây?
 \choiceTF 
  {\True  $4\sin^4x-12\sin^2x+5=0$}
   {$\cos2x=1$}
  {$4\cos^4x-4\cos^2x-3=0$} 
  {\True  $\cos2x=0$}
 \loigiai{Ta có $\cos^2x=1-\sin^2x$ nên 
  \begin{eqnarray*}
    4\sin^4x+12\cos^2x-7=0&\Leftrightarrow&4\sin^4x-12\sin^2x+5=0\\
    &\Leftrightarrow&\hoac{&\sin^2x=\dfrac{5}{2}\\&\sin^2x=\dfrac{1}{2}}\\
   &\Leftrightarrow&\cos2x=0. 
  \end{eqnarray*}
   Ta có $\sin^2x=1-\cos^2x$ nên
  \begin{eqnarray*}
   4\sin^4x+12\cos^2x-7=0&\Leftrightarrow&4\left(1-\cos^2x\right)^2+12\cos^2x-7=0\\ 
   &\Leftrightarrow&4\cos^4x+4\cos^2x-3=0. 
 \end{eqnarray*}}
\end{ex}

% nguồn: Ok_dungsai_lop11_Lan1
\begin{ex}%[1D1H6-1]%[ma:9db37d42]
 Xét phương trình $3\sin^2x-2\cos x+2=0$ $(*)$, nhận xét  nào sau đây là \textbf{sai}.
 \choiceTF 
  {\True $x=k\pi,k\in\mathbb{Z}$ là  nghiệm của phương trình $(*)$}
  {$x=k2\pi,k\in\mathbb{Z}$ là nghiệm của  phương trình $(*)$}
   {\True Nghiệm của phương trình $\sin x=0$ cũng là nghiệm của $(*)$}
  {Nghiệm của phương trình $\cos x=1$ cũng là  nghiệm của phương trình $(*)$}
 \loigiai{ 
 Ta có $3\sin^2x-2\cos  x+2=0\Leftrightarrow 3\cos^2x+2\cos x-5=0\Leftrightarrow\cos x=1\Leftrightarrow x=k2\pi,k\in\mathbb{Z}$.}
\end{ex}
